\section{Introduction}\label{sec:introduction}

A \emph{packing} of a family of matroids on a common ground set is a choice
of pairwise disjoint spanning sets, one for each matroid. A \emph{covering}
is a choice of independent sets, one for each matroid, whose union is the
ground set. The infinite matroid packing/covering conjecture asks whether
every family admits a partition of the ground set on which these two
requirements can be satisfied in the following complementary senses.

Throughout, matroids are standard infinite matroids, with the maximal
extension axiom recalled in Section~\ref{sec:matroids}. For a matroid $M$
on $E$ and $X\subseteq E$, write $M\!\upharpoonright X$ for restriction
and
\[
 M.X:=M/(E\setminus X)=(M^*\!\upharpoonright X)^*
\]
for contraction onto $X$. A set is spanning if its closure is the whole
ground set. A \emph{packing/covering partition} for $(M_i:i\in\Theta)$ is
a partition $E=P\dot\cup C$ together with sets $S_i\subseteq P$ and
$I_i\subseteq C$ such that
\begin{align*}
 S_i\cap S_j&=\varnothing &&(i\ne j),\\
 \cl_{M_i\upharpoonright P}(S_i)&=P &&(i\in\Theta),\\
 I_i&\text{ is independent in }M_i.C &&(i\in\Theta),\\
 \bigcup_{i\in\Theta}I_i&=C.
\end{align*}
Bowler and Carmesin formulated the assertion that every set-indexed family
has such a partition and established its relationship with infinite
matroid intersection~\cite[Conjecture~1.3 and Corollary~3.9]{BowlerCarmesin2015}. We give a counterexample
with two matroids.

\begin{theorem}\label{thm:main}
In ZFC, there exist a countably infinite set $E$ and two infinite matroids
$M_0,M_1$ on $E$ such that $M_i^*=M_i$ on the same labelled ground set for
$i=0,1$, and
\[
 I_0\cup I_1\ne E
 \quad\text{for every independent }I_0\text{ in }M_0
 \text{ and independent }I_1\text{ in }M_1.
\]
The pair $(M_0,M_1)$ has no packing/covering partition. Consequently, the
infinite matroid packing/covering conjecture is false.
\end{theorem}

The last conclusion needs a separate argument: failure of a covering of
$E$ alone does not rule out a mixed partition. For a self-dual pair,
however, the spanning sets on $P$ and the independent sets in the
contractions onto $C$ give disjoint spanning sets on $E$. Their complements
are an independent covering. We prove this implication in
Section~\ref{sec:matroids} before constructing the counterexample.

\subsection*{Context and the local construction}

The axioms of Bruhn, Diestel, Kriesell, Pendavingh and
Wollan~\cite{BDKPW2013} provide an infinite matroid theory with restriction,
contraction and duality. In particular, independent sets must extend to
maximal independent sets inside \emph{every} prescribed subset of the
ground set. This condition is a central part of our construction.

A matroid is \emph{uniform} if replacing any element of any basis by any
element outside that basis again gives a basis. Bowler and
Geschke~\cite[Theorem~6 and the proof of Theorem~15]{BowlerGeschke2016} developed a basis criterion for infinite
uniform matroids and constructed countable self-dual examples under
Martin's axiom for countable partially ordered sets, and hence under the
continuum hypothesis. Their criterion combines an
antichain of bases, invariance under balanced finite changes, and a basis
existence condition for intervals of subsets. Our assembly also uses the
complementary interval-recursion pattern of their Lemma~12, while
verifying the full independent-set axioms directly.

Bowler and Geschke left open whether an infinite self-dual uniform matroid
can be constructed in ZFC. Even without self-duality, ZFC existence of a
uniform matroid of infinite rank and corank is recorded as open by Gollin
and Jo\'o~\cite[Section~1]{GollinJoo2025}.

Our local construction gives a self-dual uniform matroid $Q$ on a
countable set in ZFC, with every basis and its complement infinite.
It therefore answers both of these existence questions affirmatively.
Its bases satisfy an additional ordinal inequality that obstructs
independent coverings when copies of $Q$ are placed along a doubly
infinite path. The two matroids in Theorem~\ref{thm:main} are countably
infinite direct sums of these copies. In Jo\'o's terminology, a direct sum
of uniform matroids is \emph{partitional}. His definition of uniformity is
equivalent to requiring every subset to be independent or spanning
\cite[Definition~1.2 and Proposition~2.2]{Joo2021Partitional}, the property
proved for $Q$ in Remark~\ref{rem:uniform-local}. Thus both matroids in our
pair are partitional. This distinction matters: Jo\'o proved an intersection
theorem for an arbitrary infinite matroid paired with a direct sum of
\emph{finitely many} uniform matroids, and asked whether intersection holds
for two partitional matroids on a common countable ground set
\cite[Theorem~1.4 and Question~1.5]{Joo2021Partitional}.

For matroids $M,N$ on a common ground set $E$, the pair \emph{satisfies
intersection} if there is a set $J$ independent in both matroids and a
partition $J=J_M\dot\cup J_N$ such that
\[
 \cl_M(J_M)\cup\cl_N(J_N)=E.
\]
The unrestricted infinite Matroid Intersection Conjecture asserts that
every pair of matroids on a common, possibly infinite, ground set satisfies this
property~\cite[Definition~3.5 and Section~7]{BowlerCarmesin2015}.

\begin{corollary}\label{cor:intersection}
The same pair $(M_0,M_1)$ constructed in Theorem~\ref{thm:main} consists of
partitional matroids and does not satisfy intersection. More explicitly,
for every set $J$ independent in both $M_0$ and $M_1$ and every partition
$J=J_0\dot\cup J_1$,
\[
 \cl_{M_0}(J_0)\cup\cl_{M_1}(J_1)\ne E.
\]
Consequently, the unrestricted infinite Matroid Intersection Conjecture is
false in ZFC. The pair also gives a negative ZFC answer to Jo\'o's
Question~1.5 for two partitional matroids on a common countable ground set.
\end{corollary}

\begin{proof}
Each $M_i$ is a direct sum of countably infinitely many copies of the uniform
matroid $Q$, so it is partitional. Bowler and Carmesin
\cite[Proposition~3.6]{BowlerCarmesin2015} prove that, for any matroids $M,N$
on the same ground set,
\[
 (M,N)\text{ satisfies intersection}
 \quad\Longleftrightarrow\quad
 (M,N^*)\text{ has a packing/covering partition}.
\]
Apply this equivalence with $M=M_0$ and $N=M_1$. Since $M_1^*=M_1$ on the same labelled ground set,
the right-hand side concerns the identical pair $(M_0,M_1)$ and is false by
Theorem~\ref{thm:main}. Its common ground set is countable, so it also
satisfies the hypotheses of Jo\'o's question.
\end{proof}

Bowler and Carmesin also formulated separate Covering and Packing
conjectures. Their universal equivalences imply that both conjectures are
false in ZFC; Section~\ref{sec:separate-consequences} gives the precise
formulations and deduction. This consequence concerns universal statements
about matroid families and does not identify their witnessing families
with the pair $(M_0,M_1)$.

The examples are neither finitary nor cofinitary. Here a matroid is
\emph{finitary} if a set is independent whenever all its finite subsets
are independent, and \emph{cofinitary} if its dual is finitary. Nash-Williams' original intersection conjecture concerned finitary
matroids; Aharoni and Ziv~\cite{AharoniZiv1998} gave an early published
formulation. Jo\'o proved intersection for two finitary matroids on a countable
ground set~\cite[Theorem~1.3]{Joo2021Countable}.
His subsequent packing/covering theorem allows every set-indexed family
on a countable ground set whose members are direct sums of finitary and
cofinitary matroids~\cite[Theorem~1.2]{Joo2024Packing}.
The examples here have no nonempty finitary or cofinitary direct summand,
as we show in Section~\ref{sec:counterexample}, so they lie outside this
mixed direct-sum class. Countability of the ground set is distinct from
these structural hypotheses; the present result concerns unrestricted
infinite matroids and does not refute Nash-Williams' original finitary
conjecture.

\subsection*{Proof overview}

We first construct an ultrafilter $\mathcal U$ on a countable set $D$ and
an ordinal-valued function $\rho$ on its members. Call a subset of $D$
\emph{large} if it belongs to $\mathcal U$. This auxiliary ordinal rank, distinct from matroid rank, satisfies
$\rho(X)\ge\rho(Y)$ whenever $X\subseteq Y$ are large. On two labelled copies $L,R$ of $D$, we construct
a self-dual matroid $Q$ with the following property: if $B$ is a basis
and its slice $B_L\subseteq D$ is large, then $D\setminus B_R$ is large
and
\[
 \rho(B_L)>\rho(D\setminus B_R).
\]
This is the property used at the end of the proof. To build the pair,
join consecutive integer vertices by \emph{bundles}, each a labelled
copy of $D$, and place $Q$ on the two incident bundles at each vertex.
Its $L$-coordinate points toward the central bundle between $-1$ and $0$,
and its $R$-coordinate points outward. The copies at even and odd vertices
form $M_0$ and $M_1$, respectively. In a putative independent covering, one starts with the large side of a
central bundle and follows the path outward. Each successive copy of
$Q$ strictly decreases the ordinal rank, an impossibility.

Most of the proof constructs $Q$. Partition $D$ into rapidly growing
finite blocks. The first ideal consists of subsets whose proportions
in these blocks tend to zero. Within each equivalence class modulo this
ideal, we define real-valued probes of a subset; a half-open interval
condition on their lower or upper limits selects the bases. Finite
changes shift these limits by their net cardinality, making basis
exchange immediate. A separate interpolation argument supplies bases
between suitable endpoints, including endpoints with infinite limits.
This argument is the key to maximal extension.

A second, larger ideal permits the insertion of sets that preserve the
ordinal inequalities while separating previously chosen classes.
Independent binary columns supply these separating sets simultaneously
for fewer than continuum many requirements. We then address every
interval of subsets in a transfinite recursion, inserting complementary
classes of bases when necessary. At each stage there are fewer than
continuum many earlier requirements; no regularity assumption on the
continuum is used.

Section~\ref{sec:filters} constructs the two ideals, the ultrafilter, the
rank and the simultaneous reservations. Section~\ref{sec:local-bases}
proves the interpolation properties within a single class.
Section~\ref{sec:assembly} assembles the classes and proves all the
matroid axioms, including maximal extension.
Section~\ref{sec:counterexample} places $Q$ along the path and completes
the ordinal descent argument. Section~\ref{sec:separate-consequences}
derives the separate universal Covering and Packing consequences.
